%!TeX program = xelatex
\documentclass[12pt,hyperref,a4paper,UTF8]{ctexart}
\usepackage{FzuReport}

\begin{document}

\title{ \vspace{3cm} \heiti \Huge \textbf{{编译原理课程项目\\Mandrill语言编译器实践报告}} \par
\vspace{1cm} 
\heiti \Large {\underline{Mandrill v2.0 完整实现}} \par   
\vspace{3cm} }%标题

\author{
    \vspace{0.5cm}
    \kaishu\Large 姓名\ \dlmu[9cm]{ 这里改成你的名字 } \qquad \\ %姓名 
    \vspace{0.5cm}
    \kaishu\Large  学号\ \dlmu[9cm]{ 这里改成你的学号 } \qquad  \\ %学号
    \vspace{0.5cm}
    \kaishu\Large   授课教师\ \dlmu[9cm]{ 改成你的授课教师 } %
}
\date{\today} % 默认为今天的日期，改成\today{}可以不显示日期，或者填入你期望的日期

\cover

\thispagestyle{empty} 

\newpage
\tableofcontents

\newpage

\begin{abstract}
请在这里撰写摘要，简要介绍你在编译原理课程中实现 Mandrill v2.0 语言编译器的完整过程。建议包含以下内容：

\begin{itemize}
    \item 你使用的编程语言与构建工具；
    \item 词法分析器和语法分析器采用的技术路线（按课程要求应为手写递归下降实现）；
    \item 语义分析器与编译器后端是否借助了 ANTLR；
    \item 编译器输出的目标代码形式（文本汇编指令），以及由课程组提供的模拟器执行；
    \item 测试通过情况（合法用例与语义错误用例）。
\end{itemize}

摘要应控制在 200--300 字左右，语言简洁凝练。
\end{abstract}

\section{词法分析器}

\textbf{【本章写作提示】}请按课程要求，描述你手写实现的词法分析器。以下给出建议的写作结构与参考内容，请根据你的实际实现进行补充或调整。

\subsection{整体架构与字符扫描}

请描述你的词法分析器采用的核心 IO 组件（如 \texttt{PushbackReader}、\texttt{BufferedReader} 等），以及如何逐字符读取源代码。建议说明以下技术细节：
\begin{itemize}
    \item 如何维护行号与列号（遇到换行符时的处理）；
    \item 扫描每个 Token 前如何记录起始位置；
    \item 如何处理 Windows 换行符 \texttt{\textbackslash r\textbackslash n} 或 \texttt{\textbackslash r} 的兼容性。
\end{itemize}

\subsection{多字符运算符与关键字识别}

请描述单字符运算符与多字符运算符的识别逻辑。对于多字符运算符（如 \texttt{==}、\texttt{!=}、\texttt{<=}、\texttt{>=}），请说明你如何实现向前看（look-ahead）机制。

关键字的识别建议通过查找表（如 \texttt{HashMap<String, TokenType>}）实现，请说明你的设计方案及其时间复杂度优势。

\subsection{字符串与字符常量}

请描述字符串常量与字符常量的解析逻辑，包括但不限于：
\begin{itemize}
    \item 转义字符（\texttt{\textbackslash\textbackslash}、\texttt{\textbackslash n}、\texttt{\textbackslash '} 等）的处理方式；
    \item 未闭合字符串或字符常量的错误检测与报告。
\end{itemize}

\subsection{Token 输出与错误处理}

请说明你的 Token 输出格式，以及遇到非法输入（如非法字符、未闭合字符串）时的错误处理策略（是否输出错误信息、是否安全退出等）。

词法分析器需要识别的 Token 类别可参考 \autoref{tab:tokens} 整理。你可以直接使用或修改下表。

\begin{table}[!htbp]
    \centering
    \caption{Mandrill 词法 Token 类别（示例）}
    \label{tab:tokens}
    \begin{tabular}{l|l}
    \hline
    类别 & 具体符号 \\ \hline
    关键字 & \texttt{if}, \texttt{else}, \texttt{while}, \texttt{read}, \texttt{put}, \texttt{write}, \texttt{get}, \texttt{func}, \texttt{global}, \texttt{local}, \texttt{return}, \texttt{break}, \texttt{continue} \\
    运算符 & \texttt{+}, \texttt{-}, \texttt{*}, \texttt{/}, \texttt{\%}, \texttt{<}, \texttt{<=}, \texttt{>}, \texttt{>=}, \texttt{==}, \texttt{!=}, \texttt{=} \\
    分隔符 & \texttt{(}, \texttt{)}, \texttt{[}, \texttt{]}, \texttt{\{}, \texttt{\}}, \texttt{,}, \texttt{;} \\
    字面量 & 整数、字符常量、字符串常量 \\
    标识符 & \texttt{[a-zA-Z\_][a-zA-Z0-9\_]*} \\ \hline
    \end{tabular}
\end{table}

\subsection{设计反思（可选）}

请补充你在手写词法分析器过程中的感悟。例如：哪些细节花费了比你预期更多的时间？手写实现与自动化工具相比有何优劣？

\section{语法分析器}

\textbf{【本章写作提示】}请描述你手写实现的语法分析器。以下给出建议的写作结构与参考内容。

\subsection{递归下降 Parser 的设计}

请描述你为每个非终结符编写的递归下降方法，以及如何处理 Mandrill EBNF 文法。建议重点说明：
\begin{itemize}
    \item 表达式优先级分层的设计（如何自然地消除左递归并处理优先级）；
    \item 以 \texttt{a + b * c} 为例，解释分层结构如何保证 \texttt{b * c} 被优先解析。
\end{itemize}

以下是一段产生式分层结构的 LaTeX 示例，你可以直接引用或替换为你自己的文法：
\begin{verbatim}
Expression       ::= LogicalExp
LogicalExp       ::= RelationalExp { ("==" | "!=") RelationalExp }
RelationalExp    ::= AdditiveExp { ("<" | "<=" | ">" | ">=") AdditiveExp }
AdditiveExp      ::= MultiplicativeExp { ("+" | "-") MultiplicativeExp }
MultiplicativeExp::= Primary { ("*" | "/" | "%") Primary }
\end{verbatim}

\subsection{特殊语法结构的处理}

请说明你在实现中遇到的特殊语法结构及其处理方式，例如：
\begin{itemize}
    \item 数组初始化语句 \texttt{a[] = expr}：如何通过向前看多个 Token 区分数组初始化与数组元素访问？
    \item \texttt{if-else} 语句的 dangling-else 歧义：你采用了哪种匹配策略？
    \item 空语句 \texttt{;} 的处理方式。
\end{itemize}

\subsection{错误处理策略}

请描述你的语法分析器如何处理语法错误。例如：是否先检查词法错误？遇到非预期 Token 时如何报告？是否实现了错误恢复机制（如 panic mode）？

\subsection{与 LL(1) 文法的关系（可选）}

请分析 Mandrill 文法是否满足 LL(1) 条件。对于每个非终结符的多个产生式，它们的 FIRST 集是否互不相交？你如何通过首 Token 唯一确定产生式？

\section{语义分析}

\textbf{【本章写作提示】}请描述你基于 ANTLR 实现的语义分析器。以下给出建议的写作结构与参考内容。

\subsection{符号表设计}

请描述你设计的符号表数据结构。建议说明以下三层结构及其字段：
\begin{itemize}
    \item \textbf{全局变量}：名称、类型、全局数据区编号等；
    \item \textbf{函数定义}：名称、参数列表、返回值类型、入口标签、局部变量表等；
    \item \textbf{局部变量}：名称、类型、栈帧偏移等。
\end{itemize}

请说明你是否采用两遍扫描策略（先收集签名再检查函数体），以及为什么需要这样做（如 Mandrill 支持函数前向调用）。

\subsection{类型推断规则}

Mandrill 具有隐式类型推断特性。请整理并描述你实现的类型推断规则，例如：
\begin{itemize}
    \item \texttt{a = expr} 时如何根据 \texttt{expr} 的类型推断 \texttt{a} 的类型；
    \item \texttt{a[] = expr} 时 \texttt{a} 的类型推断；
    \item \texttt{local a;} 与 \texttt{local a[];} 的区别；
    \item 未声明直接使用的变量的默认推断规则。
\end{itemize}

\subsection{类型检查}

请描述你的类型检查器实现了哪些检查规则。建议包括但不限于：
\begin{itemize}
    \item 赋值左右类型兼容性；
    \item 数组索引类型与被索引对象类型；
    \item \texttt{break}/\texttt{continue} 是否仅出现在循环体内；
    \item \texttt{return} 返回值与函数声明的一致性；
    \item 函数调用实参与形参的个数及类型匹配。
\end{itemize}

请说明错误报告机制：如何携带行列号信息？输出格式是什么？

\subsection{作用域与变量查找（可选）}

请描述 Mandrill 的作用域规则以及你的变量查找策略。例如：函数内部如何访问全局变量？局部变量、函数参数、全局变量的查找顺序是什么？

\section{Mandrill 编译器}

\textbf{【本章写作提示】}请描述你基于 ANTLR 实现的编译器。以下给出建议的写作结构与参考内容。

\subsection{中间代码与代码生成策略}

请描述你的编译器采用的代码生成策略。例如：
\begin{itemize}
    \item 是直接由 AST 生成汇编，还是经由中间表示（如三地址码 TAC）再生成汇编？
    \item 若使用了 TAC，请列出你实现的 TAC 指令种类；
    \item 临时变量的分配策略是什么？
\end{itemize}

\subsection{控制流翻译}

请描述控制流语句的翻译方案。建议画出或说明以下结构生成的汇编代码模板：
\begin{itemize}
    \item \texttt{while} 循环：条件判断、循环体、跳转回起始标签的机制；
    \item \texttt{if-else}：条件判断、then 分支、else 分支、合并标签；
    \item \texttt{break} 与 \texttt{continue}：如何通过标签栈实现目标地址回填？
\end{itemize}

请特别注意 \texttt{eval} 条件跳转的栈操作顺序：先压入 ELSE 地址，再压入 THEN 地址，最后压入条件值（cond 在栈顶）。

\subsection{函数调用的翻译}

请描述函数调用的完整汇编生成过程。建议说明：
\begin{itemize}
    \item 调用者如何压入参数；
    \item \texttt{jal} 指令的作用（保存 SP 与 PC）；
    \item 被调用者如何将参数从栈中取出并存入局部变量；
    \item 返回值如何传递；
    \item \texttt{ret} 指令如何恢复调用者上下文；
    \item 无返回语句的函数是否需要特殊处理（如默认返回 0）。
\end{itemize}

\subsection{数组与字符串的处理}

请描述编译器如何处理数组与字符串：
\begin{itemize}
    \item 数组元素访问的地址计算公式（如 $\text{Address} = \text{Base}(a) + i \times 4$）；
    \item 字符串常量的内存分配与初始化过程（\texttt{malloc}、逐字符写入、UTF-32 编码、0 结尾）；
    \item 数组赋值与字符串赋值在代码生成上的区别。
\end{itemize}

编译器生成的核心指令集可参考 \autoref{tab:instructions} 整理。

\begin{table}[!htbp]
    \centering
    \caption{Mandrill 虚拟机汇编指令集（示例）}
    \label{tab:instructions}
    \begin{tabular}{l|l|l}
    \hline
    类别 & 指令 & 说明 \\ \hline
    数据存取 & \texttt{dstore x} & 将立即数 \texttt{x} 压入操作数栈 \\
     & \texttt{dload x} & 从全局变量区读取变量 \texttt{x} 入栈 \\
     & \texttt{dlload x} & 从局部变量区（相对 SP 偏移 \texttt{x}）读取入栈 \\
     & \texttt{dlwrite x} & 将栈顶元素写入局部变量 \texttt{x} 并弹出 \\
     & \texttt{daload x} & 从栈顶地址读取数组元素入栈 \\
     & \texttt{dawrite x} & 将值写入栈顶地址指向的内存 \\
    运算 & \texttt{eval x} & 根据操作数 \texttt{x} 执行二元运算或条件跳转 \\
    跳转 & \texttt{jump x} & 将 PC 设置为 \texttt{x} \\
     & \texttt{jal x} & 链接跳转，保存上下文并调用函数 \\
     & \texttt{ret x} & 恢复调用者上下文并返回 \\
    内存 & \texttt{malloc x} & 从栈顶获取大小，申请堆内存并将地址入栈 \\
    I/O & \texttt{geti}/\texttt{getc}/\texttt{gets} & 读取整数/字符/字符串并入栈 \\
     & \texttt{puti}/\texttt{putc}/\texttt{puts} & 弹出栈顶并输出到 stdout \\ \hline
    \end{tabular}
\end{table}

\subsection{汇编模拟器执行}

请描述实验四的完整执行流程。注意：模拟器由课程组提供，你不需要实现它，但需要理解其工作原理以正确生成汇编。建议说明：
\begin{itemize}
    \item 编译器从标准输入读取源程序，输出文本汇编到标准输出；
    \item 评测机如何将汇编重定向到临时文件并由模拟器加载执行；
    \item 模拟器的文本汇编解析方式（逐行解析、注释处理、操作数格式）；
    \item 模拟器的运行时内存模型（操作数栈、全局变量区、局部变量区、调用栈帧、堆内存）；
    \item 理解模拟器对你正确生成汇编有何帮助。
\end{itemize}

\section{心得体会}

\textbf{【本章写作提示】}请结合你的实际实现过程，撰写真实的感悟与反思。以下方向供参考：

\begin{itemize}
    \item 手写词法分析器和语法分析器时，哪些细节花费了比你预期更多的时间？（如行列号维护、向前看机制、dangling-else 处理等）
    \item 从手写前端转向 ANTLR 实现语义分析与代码生成，你有何体会？工具带来了哪些便利，又有哪些原理性的东西必须手动处理？
    \item 编译器后端中，控制流跳转、函数调用约定、栈操作管理是否是难点？你是如何理清思路的？（如手绘栈状态图、单步跟踪等）
    \item 中间表示（如 TAC）是否帮助你解耦了前端与后端？
    \item 项目采用 Gradle 多模块管理，你在模块划分、依赖管理、独立测试方面有何体会？
\end{itemize}

请避免空洞的套话，尽量写出具体、真实的技术反思。

\section{总结}

\textbf{【本章写作提示】}请用 2--3 段话总结你的项目成果。建议包含：
\begin{itemize}
    \item 你完成了哪些模块的实现；
    \item 采用了哪些关键技术；
    \item 测试通过情况；
    \item 是否达到了课程实践的教学目标。
\end{itemize}

\section*{致谢}

\textbf{【本章写作提示】}请写上对你实现该项目帮助较大的同学、学长学姐、大模型或课程组教师。请保持真诚简洁。

\newpage
\appendix
\section{AI 使用情况说明}

\textbf{【必填】}根据课程学术诚信要求，请如实填写你在完成本实验过程中使用 AI 辅助的情况。若未使用任何 AI 工具，请明确写明"本次实验未使用任何 AI 工具"。

\subsection{使用声明}

\begin{table}[!htbp]
    \centering
    \begin{tabular}{|p{4cm}|p{9cm}|}
    \hline
    \textbf{项目} & \textbf{填写内容} \\
    \hline
    是否使用 AI 工具 & 是 / 否 \\
    \hline
    使用的 AI 模型 & （如：GPT-4、Claude 3.5、Kimi、Copilot、通义千问等） \\
    \hline
    使用场景 & （如：代码调试、算法设计、报告撰写、LaTeX 排版、概念解释等） \\
    \hline
    使用频率 & （如：偶尔、经常、全程辅助等） \\
    \hline
    \end{tabular}
\end{table}

\subsection{详细使用记录}

请按时间顺序或按模块列出你使用 AI 的具体情况。建议每条记录包含以下要素：
\begin{enumerate}
    \item \textbf{使用场景}：你在做什么时使用了 AI？
    \item \textbf{输入提示词（Prompt）}：你向 AI 提出的具体问题或指令（请尽量如实还原）。
    \item \textbf{AI 输出摘要}：AI 给出了什么样的回答或建议？
    \item \textbf{你的采纳与修改}：你是否直接采用了 AI 的输出？如果进行了修改，请说明原因。
\end{enumerate}

\tbox{
    \textbf{填写示例}（请删除此文本框后如实填写）：\\[0.5em]
    \textbf{场景}：编写递归下降语法分析器时遇到 dangling-else 问题。\\
    \textbf{Prompt}："如何在递归下降 parser 中处理 if-else 的 dangling-else 歧义？请给出 Java 代码示例。"\\
    \textbf{AI 输出}：建议采用贪心策略，让 else 匹配最近的未匹配 if。\\
    \textbf{采纳情况}：参考了该思路，但根据自己的文法结构重新实现了代码。
}

\vspace{1em}
\noindent\textbf{请在此处开始填写你的 AI 使用记录：}

\vspace{2em}

% 学生填写区域开始
% 请删除以下示例文本，填写你自己的记录

\textbf{记录 1}

\begin{tabular}{|p{3cm}|p{10cm}|}
\hline
场景 & \\ \hline
Prompt & \\ \hline
AI 输出摘要 & \\ \hline
采纳与修改 & \\ \hline
\end{tabular}

\vspace{1em}

\textbf{记录 2}

\begin{tabular}{|p{3cm}|p{10cm}|}
\hline
场景 & \\ \hline
Prompt & \\ \hline
AI 输出摘要 & \\ \hline
采纳与修改 & \\ \hline
\end{tabular}

\vspace{1em}

% 学生填写区域结束

\subsection{诚信承诺}

本人承诺：以上关于 AI 使用情况的说明真实、完整。本人理解并遵守课程关于学术诚信的相关规定，未将 AI 生成内容未经审阅直接作为自己的成果提交。

\vspace{2em}
\noindent 签名：\dlmu[6cm]{} \hspace{2em} 日期：\dlmu[4cm]{}

\end{document}
